\subsection{量子/时空骨架的工作边界与后续目标}\label{subsec:conclusion-quantum-spacetime-skeleton-boundary}

本文在结论层面真正回答的是如下问题：在不增设新的物理公理的前提下，统一上位逻辑是否足以纯粹导出量子结构与时空结构的共同骨架。就定理 \ref{thm:physical-spacetime-procedural-grand-chain} 与推论 \ref{cor:physical-spacetime-no-new-axioms} 所固定的范围而言，本文给出的回答是肯定的，但这种肯定必须被理解为对\emph{骨架}的肯定，而不是对全部标准物理形式的先验恢复。

更具体地说，本文已经把观察者重写为状态纤维上的局部参考系，把时间重写为同步延迟或精化深度，把空间重写为共同支撑与资源代价的比较结构，把因果重写为允许精化链诱导的前序关系；在此基础上，又把局部可读而全局受阻的现象表述为 \v{C}ech 型粘接障碍，把闭路复合 residual 表述为离散 holonomy 与离散曲率，并进一步指出：相位可被理解为局部 lift 之间粘接 residual 的可见投影，干涉可被理解为同一可见端点上的多重 local lift 无法同时规范消去的统计后果，传播上界可被理解为同步--传播比的极大不变量。于是，量子结构与时空结构不再作为互不相关的外加起点出现，而被压回到同一局部--全局失配机制在统计表示层与几何表示层上的共同显现。

与此同时，本文也刻意避免过度宣称。本文已经建立的是从统一上位逻辑到量子/时空骨架的纯推导链，而不是已经无条件恢复了标准形式的量子力学或广义相对论。平方型读出的唯一性、射影态类的 Hilbert 型表示、传播上界 \(c_*\) 在一般观察者变换下的完整协变性，以及离散曲率在细化极限中恢复连续几何对象的刚性与唯一性，仍需借助后续的表示定理、极限定理与刚性定理来完成。

因此，本文最根本的结论不是某一组现成方程，而是一种方法论上的后移原则：通常被当作起始物理公理的结构，可以被系统地重新解释为统一上位逻辑内部既有结构的表示后果。凡属量子态空间、平方型概率读法、传播上界协变性、连续曲率恢复以及场方程刚性等问题，本文一律不作为起始输入，而统一后移为表示层与极限层的目标。本文已经完成的是骨架导出；尚未完成的是表示唯一性与连续极限的闭合。

若这一程序最终闭合，则“量子”与“时空”都将不再被视为理论的起点，而只是同一抽象结构在不同尺度、不同表示层上的自然外观。
